%% Guided notes, SAT Math, Domain 2: Advanced Math (skills 2.1 to 2.3).
\newcommand{\gr}[2]{\begin{center}\begin{tikzpicture}\begin{axis}[gridaxis,#1]#2\end{axis}\end{tikzpicture}\end{center}}
\newcommand{\pt}[1]{\node[circle,fill=apblue,inner sep=1.5pt] at (axis cs:#1){};}

%%DOC 2.1 | Equivalent Expressions
\goals{\item Add, subtract and multiply polynomials.
\item Factor with a common factor, a difference of squares, a perfect square and a trinomial.
\item Use exponent rules and rational exponents, and simplify rational expressions.}
\vocab{polynomial, like terms, factor, difference of squares, perfect square trinomial, rational exponent, rational expression, excluded value.}

\subsection*{1. Operations on polynomials}
\begin{keyidea}
To subtract a polynomial, distribute the \fillin{minus sign} to every term. To multiply, multiply \fillin{each term} of the first factor by each term of the second, then combine \fillin{like terms}.
\end{keyidea}

\ex{Example 1}
Simplify $(3x^2-2x+5)-(x^2+4x-1)$ and $(2x+3)(x-4)$.
\begin{workspace}[2.4cm]
$3x^2-2x+5-x^2-4x+1=2x^2-6x+6$. \quad $(2x+3)(x-4)=2x^2-8x+3x-12=2x^2-5x-12$.
\end{workspace}

\subsection*{2. Factoring patterns}
\begin{keyidea}
\begin{tabular}{@{}ll@{}}
Difference of squares: & $a^2-b^2=$ \fillin{$(a-b)(a+b)$}\\
Perfect squares: & $a^2+2ab+b^2=$ \fillin{$(a+b)^2$}, \quad $a^2-2ab+b^2=$ \fillin{$(a-b)^2$}\\
Trinomial $x^2+bx+c$: & find two numbers with product \fillin{$c$} and sum \fillin{$b$}
\end{tabular}
\end{keyidea}

\ex{Example 2}
Factor $9x^2-25$, $\;4x^2-12x+9$ and $x^2-5x-14$.
\begin{workspace}[2.4cm]
$(3x-5)(3x+5)$. \quad $(2x-3)^2$ since $2(2x)(3)=12x$. \quad $-7$ and $2$ have product $-14$ and sum $-5$: $(x-7)(x+2)$.
\end{workspace}

\ex{Example 3 \textbullet\ Matching coefficients}
If $(ax+3)(2x-5)=6x^2+bx-15$ for all $x$, find $a$ and $b$.
\begin{workspace}[2.2cm]
Expand: $2ax^2+(6-5a)x-15$. So $2a=6$, $a=3$, and $b=6-5(3)=-9$.
\end{workspace}

\subsection*{3. Exponents and radicals}
\begin{keyidea}
$x^a\cdot x^b=$ \fillin{$x^{a+b}$}, \quad $\dfrac{x^a}{x^b}=$ \fillin{$x^{a-b}$}, \quad $(x^a)^b=$ \fillin{$x^{ab}$}, \quad $x^{-a}=\dfrac{1}{x^a}$, \quad $x^{m/n}=$ \fillin{$\sqrt[n]{x^m}$}.
\end{keyidea}

\ex{Example 4}
Simplify $(2x^3y)(3x^2y^4)$ and write $x^{3/4}$ with a radical.
\begin{workspace}[1.8cm]
$6x^5y^5$. \quad $x^{3/4}=\sqrt[4]{x^3}$ (the denominator is the index of the root).
\end{workspace}

\subsection*{4. Rational expressions}
\begin{keyidea}
Factor the numerator and the denominator, then cancel common \fillin{factors} (never terms). To add fractions, use a common \fillin{denominator}.
\end{keyidea}

\ex{Example 5}
Simplify $\dfrac{x^2-9}{x^2+5x+6}$ and $\dfrac{1}{x-1}+\dfrac{2}{x+1}$.
\begin{workspace}[2.6cm]
$\dfrac{(x-3)(x+3)}{(x+2)(x+3)}=\dfrac{x-3}{x+2}$ for $x\neq-3,-2$. \quad $\dfrac{(x+1)+2(x-1)}{(x-1)(x+1)}=\dfrac{3x-1}{x^2-1}$.
\end{workspace}

\begin{aptip}[SAT tip]
To check that two expressions are equivalent, graph both in Desmos: equivalent expressions give the same graph. Or substitute a value such as $x=2$ in each.
\end{aptip}

\begin{mistakes}
\begin{itemize}
\item $(x+3)^2\neq x^2+9$: the middle term $6x$ is missing.
\item Cancelling terms: $\dfrac{x+3}{x+5}$ does not simplify to $\dfrac35$.
\end{itemize}
\end{mistakes}

\begin{quickcheck}
1. $x^2+8x+16=$ \fillin{$(x+4)^2$}. \quad 2. $x^{1/2}\cdot x^{3/2}=$ \fillin{$x^2$}. \quad 3. $25x^2-1=$ \fillin{$(5x-1)(5x+1)$}.
\end{quickcheck}

%%DOC 2.2 | Nonlinear Equations and Systems
\goals{\item Solve quadratic equations by factoring, square roots and the quadratic formula.
\item Use the discriminant to count real solutions.
\item Solve radical, rational, absolute value and exponential equations, and linear-quadratic systems.}
\vocab{quadratic equation, zero product property, quadratic formula, discriminant, extraneous solution, absolute value, system.}

\subsection*{1. Quadratic equations}
\begin{keyidea}
For $ax^2+bx+c=0$: \[x=\frac{\fillin{$-b\pm\sqrt{b^2-4ac}$}}{\fillin{$2a$}}\]
The discriminant $D=b^2-4ac$: $D>0$ gives \fillin{two} real solutions, $D=0$ gives \fillin{one}, $D<0$ gives \fillin{none}.\\
Sum of the solutions $=-\frac ba$; product $=\frac ca$.
\end{keyidea}

\ex{Example 1 \textbullet\ Three methods}
Solve $x^2-5x-14=0$, $\;(x-3)^2=25$ and $x^2-4x-1=0$.
\begin{workspace}[2.6cm]
$(x-7)(x+2)=0$: $x=7$ or $x=-2$. \quad $x-3=\pm5$: $x=8$ or $x=-2$. \quad $x=\frac{4\pm\sqrt{16+4}}{2}=2\pm\sqrt5$.
\end{workspace}

\ex{Example 2 \textbullet\ Discriminant}
For which $k$ does $x^2+kx+9=0$ have exactly one solution? How many solutions does $2x^2-3x+5=0$ have?
\begin{workspace}[2.2cm]
$k^2-36=0$: $k=6$ or $k=-6$. \quad $D=9-40=-31<0$: no real solution.
\end{workspace}

\subsection*{2. Other equations}
\begin{keyidea}
\textbf{Radical}: isolate the root, square both sides, and \fillin{check} every answer: squaring can create \fillin{extraneous} solutions.\\
\textbf{Rational}: multiply by the common denominator; reject values that make a denominator \fillin{zero}.\\
\textbf{Absolute value}: $|A|=k$ with $k\ge0$ means $A=k$ or $A=$ \fillin{$-k$}.\\
\textbf{Exponential}: write both sides with the same \fillin{base}, then set the exponents equal.
\end{keyidea}

\ex{Example 3}
Solve $\sqrt{2x+3}=x$, $\;\dfrac{5}{x-2}=\dfrac3x$, $\;|x-4|=6$ and $2^{x+1}=32$.
\begin{workspace}[3cm]
$2x+3=x^2 \Rightarrow (x-3)(x+1)=0$; $x=-1$ fails ($\sqrt1\neq-1$), so $x=3$.\\
$5x=3(x-2) \Rightarrow x=-3$. \quad $x-4=\pm6$: $x=10$ or $x=-2$. \quad $32=2^5$: $x+1=5$, $x=4$.
\end{workspace}

\subsection*{3. Linear-quadratic systems}
\begin{keyidea}
Substitute the linear equation into the quadratic one and solve the resulting quadratic. Its number of real solutions is the number of \fillin{intersection points}.
\end{keyidea}

\ex{Example 4}
Solve $y=x^2$ and $y=4x-4$.
\gr{xmin=-1.5,xmax=4.5,ymin=-1.5,ymax=9.5,width=5.5cm,height=5.6cm}{\addplot[apblue,very thick,domain=-1.5:3.1,samples=60]{x^2};\addplot[apred,very thick,domain=0.5:3.3]{4*x-4};\pt{2,4}}
\begin{workspace}[1.8cm]
$x^2=4x-4 \Rightarrow (x-2)^2=0$: one solution, $(2,4)$. The line is tangent to the parabola.
\end{workspace}

\begin{aptip}[SAT tip]
In Desmos, graph each side of an equation as $y=\dots$ and read the intersections: they are the solutions. This is the fastest check for extraneous solutions too.
\end{aptip}

\begin{quickcheck}
1. The sum of the solutions of $x^2-7x+10=0$ is \fillin{$7$}. \quad 2. $\sqrt{x}=-3$ has \fillin{no} solution. \quad 3. $3^{x}=81$ gives $x=$ \fillin{$4$}.
\end{quickcheck}

%%DOC 2.3 | Nonlinear Functions
\goals{\item Read the vertex, intercepts and extreme value of a quadratic function in each form.
\item Write and interpret exponential growth and decay models.
\item Evaluate compositions and describe transformations of graphs.}
\vocab{vertex, axis of symmetry, standard, vertex and factored forms, exponential growth, decay, growth factor, composition, translation.}

\subsection*{1. Quadratic functions}
\begin{keyidea}
\begin{tabular}{@{}ll@{}}
Vertex form $a(x-h)^2+k$: & vertex \fillin{$(h,k)$}; minimum $k$ if $a>0$, maximum if $a<0$\\
Factored form $a(x-p)(x-q)$: & $x$-intercepts \fillin{$p$ and $q$}; vertex at $x=\frac{p+q}2$\\
Standard form $ax^2+bx+c$: & $y$-intercept \fillin{$c$}; vertex at $x=$ \fillin{$-\frac{b}{2a}$}
\end{tabular}
\end{keyidea}

\ex{Example 1}
Find the vertex of $f(x)=x^2-6x+5$ and its $x$-intercepts.
\begin{workspace}[2cm]
$x=-\frac{-6}{2}=3$, $f(3)=9-18+5=-4$: vertex $(3,-4)$. $(x-1)(x-5)=0$: intercepts $1$ and $5$.
\end{workspace}

\ex{Example 2 \textbullet\ Projectile}
A ball's height in feet is $h(t)=-16t^2+64t+5$. Find the maximum height and when it occurs.
\begin{workspace}[2cm]
$t=-\frac{64}{2(-16)}=2$ s; $h(2)=-64+128+5=69$ ft. The $5$ is the starting height.
\end{workspace}

\ex{Example 3 \textbullet\ From a graph}
\gr{xmin=-2.5,xmax=4.5,ymin=-5,ymax=6,width=5.5cm,height=5.4cm}{\addplot[apblue,very thick,domain=-2:4,samples=60]{x^2-2*x-3};\pt{-1,0}\pt{3,0}\pt{1,-4}}
Write an equation of the parabola shown.
\begin{workspace}[1.6cm]
Intercepts $-1$ and $3$: $y=a(x+1)(x-3)$. At $x=1$: $a(2)(-2)=-4$, $a=1$. So $y=(x+1)(x-3)$.
\end{workspace}

\subsection*{2. Exponential functions}
\begin{keyidea}
$f(x)=a\cdot b^{x}$: $a$ is the \fillin{initial value}; $b>1$ is growth, $0<b<1$ is decay.\\
Increase by $r\%$: $b=$ \fillin{$1+r$}; decrease by $r\%$: $b=$ \fillin{$1-r$}.\\
Multiplied by $m$ every $p$ units of time: $a\cdot m^{t/p}$.
\end{keyidea}

\ex{Example 4}
A car worth $\$20{,}000$ loses $15\%$ of its value each year. A population of $500$ bacteria triples every $4$ hours. Write both models.
\begin{workspace}[2cm]
$V(t)=20000(0.85)^t$. \quad $P(t)=500\cdot3^{t/4}$.
\end{workspace}

\subsection*{3. Composition and transformations}
\begin{keyidea}
$f(g(x))$: evaluate the \fillin{inside} function first.\\
$y=f(x-h)+k$ moves the graph of $f$ \fillin{right} $h$ and \fillin{up} $k$ (for $h,k>0$).
\end{keyidea}

\ex{Example 5}
If $f(x)=x^2+1$ and $g(x)=2x-3$, find $f(g(2))$. Describe $y=f(x+4)-1$.
\begin{workspace}[1.8cm]
$g(2)=1$, $f(1)=2$. \quad The graph of $f$ moves $4$ left and $1$ down.
\end{workspace}

\begin{mistakes}
\begin{itemize}
\item $f(x-3)$ moves the graph right, not left.
\item A $15\%$ decrease uses $0.85$, not $0.15$.
\end{itemize}
\end{mistakes}

\begin{quickcheck}
1. The minimum of $g(x)=(x-2)^2+7$ is \fillin{$7$}. \quad 2. If $f(x)=a\cdot2^x$ and $f(3)=40$, then $a=$ \fillin{$5$}.
\end{quickcheck}
